% --- Archon physics formalization source begin ---
% archon:physics
% archon:covers PhyXMiniProblems/problem_phyx_mini_0803.lean
% archon:source-report reports/phyx_mini/problem_phyx_mini_0803.source.json

\chapter{Physics problem phyx\_mini\_0803}
\label{ch:PhyXMiniProblems_problem_phyx_mini_0803}

\paragraph{Problem source.}
\#\# Physical scenario

A \textbackslash{}( 50.0 \textbackslash{}mathrm\{kg\} \textbackslash{}) woman stands on a bathroom scale while riding in the elevator as shown in figure, The elevator is initially moving downward at \textbackslash{}( 10.0 \textbackslash{}mathrm\{m/s\} \textbackslash{}); it slows to a stop with constant acceleration in a distance of \textbackslash{}( 25.0 \textbackslash{}mathrm\{m\} \textbackslash{}).

\#\# Question

While the elevator is moving downward with decreasing speed, what is the reading on the scale?

\#\# Answer choices

- A: A: 590N

- B: B: 390N

- C: C: 290N

- D: D: 690N

\#\# Figure caption (auxiliary; use the image as primary evidence)

The image shows an illustration of a woman inside an elevator. The elevator is depicted with a cable system at the top, indicating it is moving downward. Inside the elevator, the woman stands on a scale.

To the right of the elevator, there is a downward arrow labeled "Moving down with decreasing speed," signaling that the elevator's speed is reducing as it descends.

Next to this, a free-body diagram displays forces acting on the woman. The arrows point vertically, with "n" denoting the normal force upwards. Below, "w = 490 N" represents the weight acting downwards. An upward arrow labeled "a\_y" indicates upward acceleration. The Y-axis is marked, and the origin is labeled "x" at the point where the forces meet.

Overall, the illustration combines mechanical elements with physical concepts to convey the dynamics of an elevator moving downward with decreasing speed.

\#\# Dataset metadata

Dynamics; Physical Model Grounding Reasoning, Multi-Formula Reasoning

\paragraph{Recorded answer/context.}
B: B: 390N

\paragraph{Figure/image path.}
/root/proposal\_for\_physic/science-mango/phyx\_mini\_run/phyx\_data/test\_image/803.png

\paragraph{Formalization target.}
create a compiling Lean file with sorry bodies at `PhyXMiniProblems/problem\_phyx\_mini\_0803.lean`. The Lean declarations must preserve the physical quantities, dimensions or dimensional roles, figure labels, governing-law hypotheses, and final relation expressed by this problem.
Use Mathlib/Physlib names found through LeanExplore where available. If a physics API is missing, introduce faithful local abstractions rather than scalar placeholder aliases.

\begin{theorem}[Physics formalization target]
\label{thm:physics:phyx_mini_0803:target}
\lean{PhyXMiniProblems.ProblemPhyXMini0803.scaleReadingWhileDescendingAndSlowing}
\uses{def:physics:phyx-mini-0803:phyxminiproblems-problemphyxmini0803-forcemagnitudeinnewtons, def:physics:phyx-mini-0803:phyxminiproblems-problemphyxmini0803-bathroomscaleelevatorsetup, def:physics:phyx-mini-0803:phyxminiproblems-problemphyxmini0803-matchessuppliedelevatorfigure, def:physics:phyx-mini-0803:phyxminiproblems-problemphyxmini0803-matcheselevatorproblemdata, def:physics:phyx-mini-0803:phyxminiproblems-problemphyxmini0803-hasphysicalelevatorparameters, def:physics:phyx-mini-0803:phyxminiproblems-problemphyxmini0803-satisfiesidealelevatordynamicslaws}
The assigned autoformalize agent should translate this physics problem into Lean declarations in the covered file, with theorem and lemma proof bodies written as `by sorry`.
\end{theorem}
\begin{proof}
This is an autoformalization task, not a proof task. Produce statements that can later be proved by the physics prover without weakening the physical model.
\end{proof}

% --- Archon declaration topology begin ---
\section{Declaration topology}
\begin{definition}[acceleration Dimension]
  \label{def:physics:phyx-mini-0803:phyxminiproblems-problemphyxmini0803-accelerationdimension}
  \lean{PhyXMiniProblems.ProblemPhyXMini0803.accelerationDimension}
  The physical dimension of acceleration, L T⁻²
\end{definition}
\begin{proof}
  This is the declared model definition of acceleration Dimension.
\end{proof}
\begin{definition}[force Dimension]
  \label{def:physics:phyx-mini-0803:phyxminiproblems-problemphyxmini0803-forcedimension}
  \lean{PhyXMiniProblems.ProblemPhyXMini0803.forceDimension}
  The physical dimension of force, M L T⁻²
\end{definition}
\begin{proof}
  This is the declared model definition of force Dimension.
\end{proof}
\begin{definition}[Mass Quantity]
  \label{def:physics:phyx-mini-0803:phyxminiproblems-problemphyxmini0803-massquantity}
  \lean{PhyXMiniProblems.ProblemPhyXMini0803.MassQuantity}
  A nonnegative physical mass, independent of its readout unit
\end{definition}
\begin{proof}
  This is the declared model definition of Mass Quantity.
\end{proof}
\begin{definition}[Length Quantity]
  \label{def:physics:phyx-mini-0803:phyxminiproblems-problemphyxmini0803-lengthquantity}
  \lean{PhyXMiniProblems.ProblemPhyXMini0803.LengthQuantity}
  A nonnegative physical distance
\end{definition}
\begin{proof}
  This is the declared model definition of Length Quantity.
\end{proof}
\begin{definition}[Signed Length Quantity]
  \label{def:physics:phyx-mini-0803:phyxminiproblems-problemphyxmini0803-signedlengthquantity}
  \lean{PhyXMiniProblems.ProblemPhyXMini0803.SignedLengthQuantity}
  A signed vertical displacement; positive means upward
\end{definition}
\begin{proof}
  This is the declared model definition of Signed Length Quantity.
\end{proof}
\begin{definition}[Speed Magnitude Quantity]
  \label{def:physics:phyx-mini-0803:phyxminiproblems-problemphyxmini0803-speedmagnitudequantity}
  \lean{PhyXMiniProblems.ProblemPhyXMini0803.SpeedMagnitudeQuantity}
  A nonnegative physical speed magnitude
\end{definition}
\begin{proof}
  This is the declared model definition of Speed Magnitude Quantity.
\end{proof}
\begin{definition}[Signed Speed Quantity]
  \label{def:physics:phyx-mini-0803:phyxminiproblems-problemphyxmini0803-signedspeedquantity}
  \lean{PhyXMiniProblems.ProblemPhyXMini0803.SignedSpeedQuantity}
  A signed vertical velocity; positive means upward
\end{definition}
\begin{proof}
  This is the declared model definition of Signed Speed Quantity.
\end{proof}
\begin{definition}[Acceleration Magnitude Quantity]
  \label{def:physics:phyx-mini-0803:phyxminiproblems-problemphyxmini0803-accelerationmagnitudequantity}
  \lean{PhyXMiniProblems.ProblemPhyXMini0803.AccelerationMagnitudeQuantity}
  \uses{def:physics:phyx-mini-0803:phyxminiproblems-problemphyxmini0803-accelerationdimension}
  A nonnegative physical acceleration magnitude
\end{definition}
\begin{proof}
  This is the declared model definition of Acceleration Magnitude Quantity.
\end{proof}
\begin{definition}[Signed Acceleration Quantity]
  \label{def:physics:phyx-mini-0803:phyxminiproblems-problemphyxmini0803-signedaccelerationquantity}
  \lean{PhyXMiniProblems.ProblemPhyXMini0803.SignedAccelerationQuantity}
  \uses{def:physics:phyx-mini-0803:phyxminiproblems-problemphyxmini0803-accelerationdimension}
  A signed vertical acceleration; positive means upward
\end{definition}
\begin{proof}
  This is the declared model definition of Signed Acceleration Quantity.
\end{proof}
\begin{definition}[Force Magnitude Quantity]
  \label{def:physics:phyx-mini-0803:phyxminiproblems-problemphyxmini0803-forcemagnitudequantity}
  \lean{PhyXMiniProblems.ProblemPhyXMini0803.ForceMagnitudeQuantity}
  \uses{def:physics:phyx-mini-0803:phyxminiproblems-problemphyxmini0803-forcedimension}
  A nonnegative physical force magnitude
\end{definition}
\begin{proof}
  This is the declared model definition of Force Magnitude Quantity.
\end{proof}
\begin{definition}[mass Readout]
  \label{def:physics:phyx-mini-0803:phyxminiproblems-problemphyxmini0803-massreadout}
  \lean{PhyXMiniProblems.ProblemPhyXMini0803.massReadout}
  \uses{def:physics:phyx-mini-0803:phyxminiproblems-problemphyxmini0803-massquantity}
  Read mass in a selected mass unit
\end{definition}
\begin{proof}
  This is the declared model definition of mass Readout.
\end{proof}
\begin{definition}[length Readout]
  \label{def:physics:phyx-mini-0803:phyxminiproblems-problemphyxmini0803-lengthreadout}
  \lean{PhyXMiniProblems.ProblemPhyXMini0803.lengthReadout}
  \uses{def:physics:phyx-mini-0803:phyxminiproblems-problemphyxmini0803-lengthquantity}
  Read a nonnegative distance in a selected length unit
\end{definition}
\begin{proof}
  This is the declared model definition of length Readout.
\end{proof}
\begin{definition}[signed Length Readout]
  \label{def:physics:phyx-mini-0803:phyxminiproblems-problemphyxmini0803-signedlengthreadout}
  \lean{PhyXMiniProblems.ProblemPhyXMini0803.signedLengthReadout}
  \uses{def:physics:phyx-mini-0803:phyxminiproblems-problemphyxmini0803-signedlengthquantity}
  Read signed vertical displacement in a selected length unit
\end{definition}
\begin{proof}
  This is the declared model definition of signed Length Readout.
\end{proof}
\begin{definition}[speed Magnitude Readout]
  \label{def:physics:phyx-mini-0803:phyxminiproblems-problemphyxmini0803-speedmagnitudereadout}
  \lean{PhyXMiniProblems.ProblemPhyXMini0803.speedMagnitudeReadout}
  \uses{def:physics:phyx-mini-0803:phyxminiproblems-problemphyxmini0803-speedmagnitudequantity}
  Read a nonnegative speed in coherent selected length and time units
\end{definition}
\begin{proof}
  This is the declared model definition of speed Magnitude Readout.
\end{proof}
\begin{definition}[signed Speed Readout]
  \label{def:physics:phyx-mini-0803:phyxminiproblems-problemphyxmini0803-signedspeedreadout}
  \lean{PhyXMiniProblems.ProblemPhyXMini0803.signedSpeedReadout}
  \uses{def:physics:phyx-mini-0803:phyxminiproblems-problemphyxmini0803-signedspeedquantity}
  Read signed vertical velocity in coherent selected length and time units
\end{definition}
\begin{proof}
  This is the declared model definition of signed Speed Readout.
\end{proof}
\begin{definition}[acceleration Magnitude Readout]
  \label{def:physics:phyx-mini-0803:phyxminiproblems-problemphyxmini0803-accelerationmagnitudereadout}
  \lean{PhyXMiniProblems.ProblemPhyXMini0803.accelerationMagnitudeReadout}
  \uses{def:physics:phyx-mini-0803:phyxminiproblems-problemphyxmini0803-accelerationmagnitudequantity}
  Read a nonnegative acceleration in coherent selected units
\end{definition}
\begin{proof}
  This is the declared model definition of acceleration Magnitude Readout.
\end{proof}
\begin{definition}[signed Acceleration Readout]
  \label{def:physics:phyx-mini-0803:phyxminiproblems-problemphyxmini0803-signedaccelerationreadout}
  \lean{PhyXMiniProblems.ProblemPhyXMini0803.signedAccelerationReadout}
  \uses{def:physics:phyx-mini-0803:phyxminiproblems-problemphyxmini0803-signedaccelerationquantity}
  Read signed vertical acceleration in coherent selected units
\end{definition}
\begin{proof}
  This is the declared model definition of signed Acceleration Readout.
\end{proof}
\begin{definition}[force Magnitude Readout]
  \label{def:physics:phyx-mini-0803:phyxminiproblems-problemphyxmini0803-forcemagnitudereadout}
  \lean{PhyXMiniProblems.ProblemPhyXMini0803.forceMagnitudeReadout}
  \uses{def:physics:phyx-mini-0803:phyxminiproblems-problemphyxmini0803-forcemagnitudequantity}
  Read force magnitude in the coherent unit induced by selected base units
\end{definition}
\begin{proof}
  This is the declared model definition of force Magnitude Readout.
\end{proof}
\begin{definition}[mass In Kilograms]
  \label{def:physics:phyx-mini-0803:phyxminiproblems-problemphyxmini0803-massinkilograms}
  \lean{PhyXMiniProblems.ProblemPhyXMini0803.massInKilograms}
  \uses{def:physics:phyx-mini-0803:phyxminiproblems-problemphyxmini0803-massquantity, def:physics:phyx-mini-0803:phyxminiproblems-problemphyxmini0803-massreadout}
  Kilogram readout used for the woman's stated mass
\end{definition}
\begin{proof}
  This is the declared model definition of mass In Kilograms.
\end{proof}
\begin{definition}[length In Meters]
  \label{def:physics:phyx-mini-0803:phyxminiproblems-problemphyxmini0803-lengthinmeters}
  \lean{PhyXMiniProblems.ProblemPhyXMini0803.lengthInMeters}
  \uses{def:physics:phyx-mini-0803:phyxminiproblems-problemphyxmini0803-lengthquantity, def:physics:phyx-mini-0803:phyxminiproblems-problemphyxmini0803-lengthreadout}
  Meter readout used for the stated stopping distance
\end{definition}
\begin{proof}
  This is the declared model definition of length In Meters.
\end{proof}
\begin{definition}[signed Length In Meters]
  \label{def:physics:phyx-mini-0803:phyxminiproblems-problemphyxmini0803-signedlengthinmeters}
  \lean{PhyXMiniProblems.ProblemPhyXMini0803.signedLengthInMeters}
  \uses{def:physics:phyx-mini-0803:phyxminiproblems-problemphyxmini0803-signedlengthquantity, def:physics:phyx-mini-0803:phyxminiproblems-problemphyxmini0803-signedlengthreadout}
  Upward-positive meter readout of the stopping displacement
\end{definition}
\begin{proof}
  This is the declared model definition of signed Length In Meters.
\end{proof}
\begin{definition}[speed Magnitude In Meters Per Second]
  \label{def:physics:phyx-mini-0803:phyxminiproblems-problemphyxmini0803-speedmagnitudeinmeterspersecond}
  \lean{PhyXMiniProblems.ProblemPhyXMini0803.speedMagnitudeInMetersPerSecond}
  \uses{def:physics:phyx-mini-0803:phyxminiproblems-problemphyxmini0803-speedmagnitudequantity, def:physics:phyx-mini-0803:phyxminiproblems-problemphyxmini0803-speedmagnitudereadout}
  Meter-per-second readout of a speed magnitude
\end{definition}
\begin{proof}
  This is the declared model definition of speed Magnitude In Meters Per Second.
\end{proof}
\begin{definition}[signed Speed In Meters Per Second]
  \label{def:physics:phyx-mini-0803:phyxminiproblems-problemphyxmini0803-signedspeedinmeterspersecond}
  \lean{PhyXMiniProblems.ProblemPhyXMini0803.signedSpeedInMetersPerSecond}
  \uses{def:physics:phyx-mini-0803:phyxminiproblems-problemphyxmini0803-signedspeedquantity, def:physics:phyx-mini-0803:phyxminiproblems-problemphyxmini0803-signedspeedreadout}
  Upward-positive meter-per-second readout of vertical velocity
\end{definition}
\begin{proof}
  This is the declared model definition of signed Speed In Meters Per Second.
\end{proof}
\begin{definition}[acceleration Magnitude In Meters Per Second Squared]
  \label{def:physics:phyx-mini-0803:phyxminiproblems-problemphyxmini0803-accelerationmagnitudeinmeterspersecondsquared}
  \lean{PhyXMiniProblems.ProblemPhyXMini0803.accelerationMagnitudeInMetersPerSecondSquared}
  \uses{def:physics:phyx-mini-0803:phyxminiproblems-problemphyxmini0803-accelerationmagnitudequantity, def:physics:phyx-mini-0803:phyxminiproblems-problemphyxmini0803-accelerationmagnitudereadout}
  Meter-per-second-squared readout of an acceleration magnitude
\end{definition}
\begin{proof}
  This is the declared model definition of acceleration Magnitude In Meters Per Second Squared.
\end{proof}
\begin{definition}[signed Acceleration In Meters Per Second Squared]
  \label{def:physics:phyx-mini-0803:phyxminiproblems-problemphyxmini0803-signedaccelerationinmeterspersecondsquared}
  \lean{PhyXMiniProblems.ProblemPhyXMini0803.signedAccelerationInMetersPerSecondSquared}
  \uses{def:physics:phyx-mini-0803:phyxminiproblems-problemphyxmini0803-signedaccelerationquantity, def:physics:phyx-mini-0803:phyxminiproblems-problemphyxmini0803-signedaccelerationreadout}
  Upward-positive meter-per-second-squared readout of acceleration
\end{definition}
\begin{proof}
  This is the declared model definition of signed Acceleration In Meters Per Second Squared.
\end{proof}
\begin{definition}[force Magnitude In Newtons]
  \label{def:physics:phyx-mini-0803:phyxminiproblems-problemphyxmini0803-forcemagnitudeinnewtons}
  \lean{PhyXMiniProblems.ProblemPhyXMini0803.forceMagnitudeInNewtons}
  \uses{def:physics:phyx-mini-0803:phyxminiproblems-problemphyxmini0803-forcemagnitudequantity, def:physics:phyx-mini-0803:phyxminiproblems-problemphyxmini0803-forcemagnitudereadout}
  Newton readout of a physical force magnitude
\end{definition}
\begin{proof}
  This is the declared model definition of force Magnitude In Newtons.
\end{proof}
\begin{definition}[Vertical Direction]
  \label{def:physics:phyx-mini-0803:phyxminiproblems-problemphyxmini0803-verticaldirection}
  \lean{PhyXMiniProblems.ProblemPhyXMini0803.VerticalDirection}
  Vertical directions occurring in the motion arrow and free-body diagram
\end{definition}
\begin{proof}
  This is the declared model definition of Vertical Direction.
\end{proof}
\begin{definition}[Speed Trend]
  \label{def:physics:phyx-mini-0803:phyxminiproblems-problemphyxmini0803-speedtrend}
  \lean{PhyXMiniProblems.ProblemPhyXMini0803.SpeedTrend}
  The trend described by “moving down with decreasing speed.”
\end{definition}
\begin{proof}
  This is the declared model definition of Speed Trend.
\end{proof}
\begin{definition}[Figure Text Label]
  \label{def:physics:phyx-mini-0803:phyxminiproblems-problemphyxmini0803-figuretextlabel}
  \lean{PhyXMiniProblems.ProblemPhyXMini0803.FigureTextLabel}
  Literal labels visible in the supplied raster
\end{definition}
\begin{proof}
  This is the declared model definition of Figure Text Label.
\end{proof}
\begin{definition}[Figure Vector Label]
  \label{def:physics:phyx-mini-0803:phyxminiproblems-problemphyxmini0803-figurevectorlabel}
  \lean{PhyXMiniProblems.ProblemPhyXMini0803.FigureVectorLabel}
  The three directed arrows in the woman's free-body/acceleration diagram
\end{definition}
\begin{proof}
  This is the declared model definition of Figure Vector Label.
\end{proof}
\begin{definition}[Elevator Figure]
  \label{def:physics:phyx-mini-0803:phyxminiproblems-problemphyxmini0803-elevatorfigure}
  \lean{PhyXMiniProblems.ProblemPhyXMini0803.ElevatorFigure}
  \uses{def:physics:phyx-mini-0803:phyxminiproblems-problemphyxmini0803-verticaldirection, def:physics:phyx-mini-0803:phyxminiproblems-problemphyxmini0803-speedtrend, def:physics:phyx-mini-0803:phyxminiproblems-problemphyxmini0803-figuretextlabel, def:physics:phyx-mini-0803:phyxminiproblems-problemphyxmini0803-figurevectorlabel}
  Qualitative and numerical information read directly from the supplied figure. The normal-force arrow has no numerical label, so this structure does not contain the requested scale reading
\end{definition}
\begin{proof}
  This is the declared model definition of Elevator Figure.
\end{proof}
\begin{definition}[Elevator Motion Regime]
  \label{def:physics:phyx-mini-0803:phyxminiproblems-problemphyxmini0803-elevatormotionregime}
  \lean{PhyXMiniProblems.ProblemPhyXMini0803.ElevatorMotionRegime}
  Motion regime asked about in the problem
\end{definition}
\begin{proof}
  This is the declared model definition of Elevator Motion Regime.
\end{proof}
\begin{definition}[Bathroom Scale Elevator Setup]
  \label{def:physics:phyx-mini-0803:phyxminiproblems-problemphyxmini0803-bathroomscaleelevatorsetup}
  \lean{PhyXMiniProblems.ProblemPhyXMini0803.BathroomScaleElevatorSetup}
  \uses{def:physics:phyx-mini-0803:phyxminiproblems-problemphyxmini0803-massquantity, def:physics:phyx-mini-0803:phyxminiproblems-problemphyxmini0803-lengthquantity, def:physics:phyx-mini-0803:phyxminiproblems-problemphyxmini0803-signedlengthquantity, def:physics:phyx-mini-0803:phyxminiproblems-problemphyxmini0803-speedmagnitudequantity, def:physics:phyx-mini-0803:phyxminiproblems-problemphyxmini0803-signedspeedquantity, def:physics:phyx-mini-0803:phyxminiproblems-problemphyxmini0803-accelerationmagnitudequantity, def:physics:phyx-mini-0803:phyxminiproblems-problemphyxmini0803-signedaccelerationquantity, def:physics:phyx-mini-0803:phyxminiproblems-problemphyxmini0803-forcemagnitudequantity, def:physics:phyx-mini-0803:phyxminiproblems-problemphyxmini0803-elevatorfigure, def:physics:phyx-mini-0803:phyxminiproblems-problemphyxmini0803-elevatormotionregime}
  Independent physical quantities in the scenario. In particular, elevatorAcceleration, normalForceN, and scaleReading are unknowns; none is defined from a displayed answer choice
\end{definition}
\begin{proof}
  This is the declared model definition of Bathroom Scale Elevator Setup.
\end{proof}
\begin{definition}[Matches Supplied Elevator Figure]
  \label{def:physics:phyx-mini-0803:phyxminiproblems-problemphyxmini0803-matchessuppliedelevatorfigure}
  \lean{PhyXMiniProblems.ProblemPhyXMini0803.MatchesSuppliedElevatorFigure}
  \uses{def:physics:phyx-mini-0803:phyxminiproblems-problemphyxmini0803-forcemagnitudeinnewtons, def:physics:phyx-mini-0803:phyxminiproblems-problemphyxmini0803-bathroomscaleelevatorsetup}
  Primary-image evidence: the motion arrow is downward, speed is decreasing, the Y axis and n/a\_y arrows are upward, and w = 490 N points downward
\end{definition}
\begin{proof}
  This is the declared model definition of Matches Supplied Elevator Figure.
\end{proof}
\begin{definition}[Matches Elevator Problem Data]
  \label{def:physics:phyx-mini-0803:phyxminiproblems-problemphyxmini0803-matcheselevatorproblemdata}
  \lean{PhyXMiniProblems.ProblemPhyXMini0803.MatchesElevatorProblemData}
  \uses{def:physics:phyx-mini-0803:phyxminiproblems-problemphyxmini0803-massinkilograms, def:physics:phyx-mini-0803:phyxminiproblems-problemphyxmini0803-lengthinmeters, def:physics:phyx-mini-0803:phyxminiproblems-problemphyxmini0803-signedlengthinmeters, def:physics:phyx-mini-0803:phyxminiproblems-problemphyxmini0803-speedmagnitudeinmeterspersecond, def:physics:phyx-mini-0803:phyxminiproblems-problemphyxmini0803-signedspeedinmeterspersecond, def:physics:phyx-mini-0803:phyxminiproblems-problemphyxmini0803-accelerationmagnitudeinmeterspersecondsquared, def:physics:phyx-mini-0803:phyxminiproblems-problemphyxmini0803-bathroomscaleelevatorsetup}
  Numerical and qualitative data from the problem prose. Signed quantities use the upward-positive convention shown in the figure. No acceleration, normal force, or scale-reading value is assumed here
\end{definition}
\begin{proof}
  This is the declared model definition of Matches Elevator Problem Data.
\end{proof}
\begin{definition}[Has Physical Elevator Parameters]
  \label{def:physics:phyx-mini-0803:phyxminiproblems-problemphyxmini0803-hasphysicalelevatorparameters}
  \lean{PhyXMiniProblems.ProblemPhyXMini0803.HasPhysicalElevatorParameters}
  \uses{def:physics:phyx-mini-0803:phyxminiproblems-problemphyxmini0803-massinkilograms, def:physics:phyx-mini-0803:phyxminiproblems-problemphyxmini0803-lengthinmeters, def:physics:phyx-mini-0803:phyxminiproblems-problemphyxmini0803-speedmagnitudeinmeterspersecond, def:physics:phyx-mini-0803:phyxminiproblems-problemphyxmini0803-accelerationmagnitudeinmeterspersecondsquared, def:physics:phyx-mini-0803:phyxminiproblems-problemphyxmini0803-bathroomscaleelevatorsetup}
  Positivity and nondegeneracy conditions on the independent magnitudes
\end{definition}
\begin{proof}
  This is the declared model definition of Has Physical Elevator Parameters.
\end{proof}
\begin{definition}[Satisfies Ideal Elevator Dynamics Laws]
  \label{def:physics:phyx-mini-0803:phyxminiproblems-problemphyxmini0803-satisfiesidealelevatordynamicslaws}
  \lean{PhyXMiniProblems.ProblemPhyXMini0803.SatisfiesIdealElevatorDynamicsLaws}
  \uses{def:physics:phyx-mini-0803:phyxminiproblems-problemphyxmini0803-massreadout, def:physics:phyx-mini-0803:phyxminiproblems-problemphyxmini0803-signedlengthreadout, def:physics:phyx-mini-0803:phyxminiproblems-problemphyxmini0803-signedspeedreadout, def:physics:phyx-mini-0803:phyxminiproblems-problemphyxmini0803-accelerationmagnitudereadout, def:physics:phyx-mini-0803:phyxminiproblems-problemphyxmini0803-signedaccelerationreadout, def:physics:phyx-mini-0803:phyxminiproblems-problemphyxmini0803-forcemagnitudereadout, def:physics:phyx-mini-0803:phyxminiproblems-problemphyxmini0803-bathroomscaleelevatorsetup}
  Governing laws for the idealized motion and scale: constant-acceleration kinematics along the upward-positive vertical axis; the woman shares the elevator's acceleration while remaining on the scale; w = m g for the woman's weight magnitude; n - w = m a\_y, Newton's second law for the woman; an ideal bathroom scale reports the normal-force magnitude. All equations hold in arbitrary coherent base units. These laws constrain the unknown scale reading but do not assign it the requested numerical value
\end{definition}
\begin{proof}
  This is the declared model definition of Satisfies Ideal Elevator Dynamics Laws.
\end{proof}
\begin{definition}[Answer Choice]
  \label{def:physics:phyx-mini-0803:phyxminiproblems-problemphyxmini0803-answerchoice}
  \lean{PhyXMiniProblems.ProblemPhyXMini0803.AnswerChoice}
  Labels of the four displayed scale-reading choices
\end{definition}
\begin{proof}
  This is the declared model definition of Answer Choice.
\end{proof}
\begin{definition}[displayed Scale Reading In Newtons]
  \label{def:physics:phyx-mini-0803:phyxminiproblems-problemphyxmini0803-displayedscalereadinginnewtons}
  \lean{PhyXMiniProblems.ProblemPhyXMini0803.displayedScaleReadingInNewtons}
  \uses{def:physics:phyx-mini-0803:phyxminiproblems-problemphyxmini0803-answerchoice}
  Every displayed choice, read in newtons
\end{definition}
\begin{proof}
  This is the declared model definition of displayed Scale Reading In Newtons.
\end{proof}
\begin{definition}[recorded Dataset Answer]
  \label{def:physics:phyx-mini-0803:phyxminiproblems-problemphyxmini0803-recordeddatasetanswer}
  \lean{PhyXMiniProblems.ProblemPhyXMini0803.recordedDatasetAnswer}
  \uses{def:physics:phyx-mini-0803:phyxminiproblems-problemphyxmini0803-answerchoice}
  Dataset answer metadata, deliberately not used as a theorem premise
\end{definition}
\begin{proof}
  This is the declared model definition of recorded Dataset Answer.
\end{proof}
\begin{lemma}[elevator Acceleration From Stopping Data]
  \label{lem:physics:phyx-mini-0803:phyxminiproblems-problemphyxmini0803-elevatoraccelerationfromstoppingdata}
  \lean{PhyXMiniProblems.ProblemPhyXMini0803.elevatorAccelerationFromStoppingData}
  \uses{def:physics:phyx-mini-0803:phyxminiproblems-problemphyxmini0803-signedaccelerationinmeterspersecondsquared, def:physics:phyx-mini-0803:phyxminiproblems-problemphyxmini0803-bathroomscaleelevatorsetup, def:physics:phyx-mini-0803:phyxminiproblems-problemphyxmini0803-matcheselevatorproblemdata, def:physics:phyx-mini-0803:phyxminiproblems-problemphyxmini0803-satisfiesidealelevatordynamicslaws}
  The stopping data determine an upward elevator acceleration of 2 m/s²
\end{lemma}
\begin{proof}
  The conclusion follows from the governing relations and figure readouts named in the statement of elevator Acceleration From Stopping Data.
\end{proof}
% --- Archon declaration topology end ---
% --- Archon physics formalization source end ---
